%%==================================================
%% chapter02.tex for Doctoral Thesis
%% Chapter 2: Theoretical Background
%% Structure: LLM -> Scene generation -> Asset reconstruction
%%==================================================

% !TeX root = ../main.tex

\chapter{理论基础}
\enchapter{Theoretical Background}
\label{chap:theory}

\section{大语言模型基础}
\ensection{Theoretical Foundations of Large Language Models}
\label{sec:llm_foundation}

\subsection{Transformer网络}
\ensubsection{Transformer}
\label{sec:transformer}

Transformer\cite{vaswani2017attention} 是一种基于注意力机制的序列建模网络结构，与循环神经网络（Recurrent Neural Network\cite{elman1990finding}, RNN）或卷积神经网络（Convolutional Neural Network\cite{lecun1998gradient}, CNN）不同，Transformer 不再依赖递归的时序状态传递，而是通过自注意力机制直接建模序列中任意位置之间的依赖关系。因此，该结构能够在保持全局建模能力的同时实现更高程度的并行计算，已成为大语言模型的核心网络基础。对于本文后续涉及的大语言模型推理、提示词引导以及结构化文本生成任务而言，Transformer 提供了最基本的表示学习与序列生成机制。

Transformer 采用编码器（Encoder）与解码器（Decoder）组成的序列到序列架构。设输入序列为
\begin{equation}
X = (\text{x}_1,\text{x}_2,\dots,\text{x}_n),
\end{equation}
其中，$\text{x}_i$ 表示第 $i$ 个位置的输入嵌入向量。编码器的作用是将输入序列映射为一组上下文相关的隐藏表示，解码器则在已生成历史信息与编码结果的条件下，自回归地预测输出序列。虽然当前大语言模型中广泛采用的是仅保留解码器的结构变体，但其核心计算机制仍然建立在 Transformer 的注意力模块与前馈网络之上。其整体网络结构如图~\ref{fig:transformer_architecture} 所示，编码器由多头自注意力与前馈网络堆叠组成，解码器在此基础上增加掩码多头注意力以及编码器-解码器注意力模块。

\begin{figure}[!t]
    \centering
    \includegraphics[width=0.8\textwidth, height=14cm]{figures/ch2/transformer_architecture}
    \caption{Transformer 网络架构图\cite{vaswani2017attention}。}
    \label{fig:transformer_architecture}
\end{figure}

Transformer 的核心贡献在于自注意力（Self-Attention）机制。对于输入特征矩阵 $X$，网络首先通过线性变换得到查询（Query）、键（Key）和值（Value）：
\begin{equation}
Q = XW^Q,\qquad
K = XW^K,\qquad
V = XW^V,
\end{equation}
其中，$W^Q$、$W^K$ 和 $W^V$ 为可学习参数矩阵。自注意力的基本思想是：对序列中每一个位置，都根据其查询向量与所有位置键向量之间的相似性，计算该位置对其他位置的关注程度，再对对应的值向量进行加权聚合，从而获得包含全局上下文信息的新表示。

具体而言，Transformer 采用缩放点积注意力（Scaled Dot-Product Attention）作为基本注意力形式，其定义为
\begin{equation}
\mathrm{Attention}(Q,K,V)
=
\mathrm{softmax}\left(\frac{QK^\top}{\sqrt{d_k}}\right)V,
\label{eq:transformer_attention}
\end{equation}
其中，$d_k$ 表示键向量的维度。式中，$QK^\top$ 用于衡量不同位置之间的相关性，$\sqrt{d_k}$ 的缩放操作用于抑制高维内积过大所带来的梯度不稳定问题，而 softmax 则将相关性归一化为注意力权重。由此可见，注意力机制本质上实现了一种基于内容相关性的动态信息路由，使模型能够直接建立长距离依赖关系。

为了增强表示能力，Transformer 并不只计算单个注意力映射，而是采用多头注意力（Multi-Head Attention）机制。其基本形式可写为
\begin{equation}
\mathrm{head}_i
=
\mathrm{Attention}(QW_i^Q,KW_i^K,VW_i^V),
\end{equation}
\begin{equation}
\mathrm{MultiHead}(Q,K,V)
=
\mathrm{Concat}(\mathrm{head}_1,\dots,\mathrm{head}_h)W^O,
\label{eq:transformer_multihead}
\end{equation}
其中，$h$ 表示注意力头的数量，$W_i^Q$、$W_i^K$、$W_i^V$ 以及 $W^O$ 均为可学习参数。多头机制的作用在于将输入特征投影到多个不同的子空间，并在各子空间中分别学习依赖关系，从而使模型能够同时关注不同尺度、不同类型的语义关联。相比单头注意力，这种设计能够显著提升模型对复杂结构信息的建模能力。

在每个注意力子层之后，Transformer 还引入逐位置前馈网络（Position-wise Feed-Forward Network\cite{vaswani2017attention}, FFN）以增强非线性表达能力。其形式通常写为
\begin{equation}
\mathrm{FFN}(\text{x})
=
W_2\,\sigma(W_1\text{x}+b_1)+b_2,
\label{eq:transformer_ffn}
\end{equation}
其中，$W_1$、$W_2$、$b_1$ 和 $b_2$ 为网络参数，$\sigma(\cdot)$ 表示非线性激活函数。该模块对序列中每个位置独立地进行相同变换，因此不会显式改变不同位置之间的关系，但能够在注意力完成信息聚合之后进一步提升特征表征能力。

为了稳定深层网络训练，Transformer 在注意力层与前馈层外都使用残差连接（Residual Connection）和层归一化（Layer Normalization）。若记某一子层为 $\text{F}$，则其输出一般写为
\begin{equation}
\text{H}' = \mathrm{LayerNorm}(\text{H}+\text{F}(\text{H})).
\label{eq:transformer_residual}
\end{equation}
这种结构能够缓解深层网络中的梯度消失与训练不稳定问题，使模型在堆叠较多层数时仍能保持良好的优化性质。因此，Transformer 可以通过不断堆叠“多头注意力 + 前馈网络”的基本模块，逐层提取更加抽象的上下文特征。

由于 Transformer 本身不包含循环结构，也不具备卷积网络中隐含的局部顺序归纳偏置，因此必须显式引入位置信息。为此，采用位置编码（Positional Encoding）来表示序列中各元素的顺序，其正弦—余弦形式定义为
\begin{equation}
\mathrm{PE}(pos,2i)=\sin\left(\frac{pos}{10000^{2i/d_{\mathrm{model}}}}\right),
\end{equation}
\begin{equation}
\mathrm{PE}(pos,2i+1)=\cos\left(\frac{pos}{10000^{2i/d_{\mathrm{model}}}}\right),
\label{eq:transformer_pe}
\end{equation}
其中，$pos$ 表示位置索引，$i$ 表示维度索引，$d_{\mathrm{model}}$ 表示特征维度。位置编码与词嵌入逐元素相加后共同作为网络输入，从而使模型在注意力计算时不仅能够感知内容相关性，也能够区分不同元素的先后顺序。

在解码阶段，为保证自回归生成的因果性，Transformer 还会在解码器自注意力中引入掩码机制（Masked Self-Attention），使当前位置只能访问当前位置及其之前的历史信息，而不能看到未来位置的内容。这样，模型便能够按照从左到右的方式逐步生成输出序列。对于大语言模型而言，这一机制决定了模型能够在给定前文条件的基础上，持续预测下一个 token，并最终完成文本续写、问答或结构化生成等任务。


\subsection{提示工程}
\ensubsection{Prompt Engineering}
\label{sec:prompt_engineering}

提示工程（Prompt Engineering）是指在不更新大语言模型参数的前提下，通过设计任务表明、上下文示例与输出约束，引导大语言模型执行特定任务的算法。在结构化生成任务中，提示工程的核心作用在于利用模型在大规模预训练阶段已经习得的语言模式与任务规律，使其在当前上下文中完成对任务目标的理解、对输出结果结构的遵循。


\noindent\textbf{（1）上下文学习}

上下文学习（In-context Learning, ICL）是大语言模型的一项重要能力，其基本思想是在推理阶段仅通过输入上下文中的任务描述与少量示例，使模型在当前上下文内完成任务适配。Brown 等人的工作\cite{brown2020gpt3} 将这一过程解释为一种发生在前向传播过程中的上下文内学习机制，即模型在预训练中已经形成了对多类任务模式的统计表征，在推理时只需通过提示显式指定当前任务形式，模型便能够利用已有参数完成对应的输出生成。按照提示中提供信息的多少，上下文学习通常可分为零样本、单样本和少样本三种形式：零样本仅给出自然语言任务表明，单样本提供一个示例，少样本则给出多个输入输出示例。三者的差异主要体现在提示中提供的任务先验强弱不同：零样本仅依赖自然语言指令激活模型已有知识，适用于任务形式较清晰、输出要求较直接的情形。单样本通过一个示例补充任务接口与输出格式信息，但由于示例数量有限，往往只能提供较弱的模式约束。少样本则通过多个示例更充分地展示输入输出对应关系与结构模式，通常更有利于模型识别任务规律，但也会占用更多上下文长度，并增加词元（Token）成本。

上下文学习的有效性来自两个方面。首先是自然语言的提示能够显式给模型提示出任务目标，从而缩小可能的输出空间。其次是，示例不仅提供语义上的参考，还隐含规定了输入与输出之间的映射关系、字段组织方式以及结果表达风格。对于需要生成结构化结果的任务而言，这一点尤为关键。由于自然语言输入往往具有表达灵活、信息分布分散和格式不统一等特点，而目标输出通常要求字段明确、层次清晰、格式固定，因此通过在提示中给出少量高质量示例，可以使模型在上下文中对“应当输出什么”和“应当如何输出”形成更稳定的条件约束。这意味着，上下文学习通过示例将任务的目标形式、局部规则与整体结构通过提示词写进当前上下文，从而引导模型完成从自然语言输入到结构化输出的映射。


\noindent\textbf{（2）思维链}

对于涉及多步推理、条件组合或复杂语义分解的任务，直接要求模型一次性给出最终答案，往往容易出现步骤遗漏、局部矛盾或推理跳跃等问题。为此，思维链（Chain-of-Thought\cite{wei2022chainofthought}, CoT）提示的核心是提出在提示词中不仅需要明确指出输入与最终输出，还需要详细展示中间推理过程。从形式上看，CoT 提示将传统的“输入—输出”示例扩展为“输入—推理链—输出”三元结构，其中推理链是通向最终结果的一系列中间自然语言步骤。其本质在于将原本隐含在模型内部的一步式生成过程展开为可显式表述的分步推理过程，使模型能够先完成问题分解与局部判断，再逐步汇总得到最终答案。

思维链的作用主要体现在三个方面。首先，它能够将复杂问题拆解为若干更容易处理的中间子问题，从而降低模型直接生成最终结果时面临的组合难度。其次，中间推理步骤为模型提供了更长的推理轨迹，使生成过程不再局限于最终答案本身，而是能够围绕中间状态逐步展开，这在多步算术、常识推理和符号推理任务中尤为重要。最后，显式推理链还提供了一定程度的可解释性，使模型的生成结果不只是一个终值，还伴随着可供检查的推理路径。CoT 论文的实验表明，这种算法在大规模语言模型上能够显著提升复杂推理任务的性能，其能力并非简单来自答案模板的变化，而是来自中间推理步骤对问题分解与推理组织的促进作用。

在 CoT 的基础上，Zero-shot-CoT\cite{kojima2022large} 进一步表明显式分步推理并不一定要求人工构造推理样例。该研究发现，只需在答案前加入诸如“Let’s think step by step”这样的简单触发语句，模型便能够在零样本条件下先生成中间推理步骤，再给出最终答案，并在多个算术、符号和逻辑推理基准上显著优于普通零样本样例的提示。

% \begin{figure}[htbp]
%   \centering
%   \includegraphics[width=0.95\linewidth]{figures/ch2/zero_shot_cot}
%   \caption{标准提示、Few-shot-CoT 提示与Zero-shot-CoT 提示的对比图。对于同一问题，传统 Few-shot 与 Zero-shot 提示均倾向于直接生成最终答案，容易出现错误。而在提示中显式加入中间推理过程，或通过“Let’s think step by step”诱导模型进行分步推理后，模型能够生成更合理的推理链并得到正确结果。图片来源于 Zero-shot-CoT \cite{kojima2022large}。}
%   \label{fig:zero_shot_cot}
% \end{figure}

%如图~\ref{fig:zero_shot_cot} 所示，
在传统提示方式下，模型往往直接从问题映射到答案，缺少对中间状态的显式建模，因此在涉及多步推导时容易出现错误。Few-shot-CoT 通过示例中的推理链显式规定了求解路径，而 Zero-shot-CoT 则表明，即使没有人工示例，简单的触发语句也能够将模型从直接作答的模式切换为先推理、后作答的模式。这表明思维链算法的关键并不在于具体模板本身，而在于通过中间自然语言步骤重组模型的生成过程，使其能够围绕若干中间状态逐步逼近最终答案。

\section{遗传算法}
\ensection{Genetic Algorithm}
\label{sec:genetic_algorithm}

遗传算法（Genetic Algorithm\cite{holland1992adaptation}, GA）借鉴了自然界优胜劣汰的演化逻辑，是一种基于种群演进的启发式搜索与优化算法。其核心思想是将待优化问题的候选解表示为可遗传的编码个体，在种群层面通过模拟自然选择、交叉与变异等生物遗传机制，驱动种群整体在迭代中持续进化，最终逼近全局最优解。相比于依赖梯度信息的优化算法，遗传算法不要求目标函数可导，也不依赖精确的局部结构信息，因此特别适用于解空间复杂、目标函数非凸、约束关系复杂或难以显式建模的优化问题。

遗传算法首先需要建立问题解到染色体编码之间的映射。对于给定优化问题，可将一个候选解记为个体
\begin{equation}
\text{x} = (x_1, x_2, \dots, x_L)
\label{eq:ga_individual}
\end{equation}
其中各维变量经过适当编码后构成长度为 $L$ 的染色体。经典遗传算法通常采用二进制串表示，即
\begin{equation}
\text{x} \in \{0,1\}^{L}
\label{eq:ga_binary_encoding}
\end{equation}
每一个基因位对应解空间中的一个影响因子。多个个体共同组成种群
\begin{equation}
P^{(t)}=\{\text{x}_1^{(t)}, \text{x}_2^{(t)}, \dots, \text{x}_N^{(t)}\}
\label{eq:ga_population}
\end{equation}
其中 $t$ 表示当前迭代代数，$N$ 表示种群规模。这种以种群为单位的机制，使得遗传算法在整个候选解集合上并行搜索，因此具有较强的全局探索能力，能够有效避免陷入局部最优，适合于复杂的组合优化问题。

为了比较不同个体的优劣，遗传算法为每个个体定义适应度函数
\begin{equation}
F(\text{x}) \ge 0
\label{eq:ga_fitness}
\end{equation}
其数值用于衡量该个体对当前优化目标的匹配程度。若原问题为最大化问题，则适应度函数通常可直接由目标函数构造。若原问题为最小化问题，则可通过单调变换将其转化为适应度越大表示个体越优的形式。在第 $t$ 代种群中，个体 $\text{x}_i^{(t)}$ 的相对适应度决定了其被保留与参与繁殖的概率，因此适应度评估构成了整个遗传搜索过程的驱动力。对于标准比例选择，个体被选中的概率可写为
\begin{equation}
p_i^{(t)}=\frac{F(\text{x}_i^{(t)})}{\sum_{j=1}^{N}F(\text{x}_j^{(t)})}
\label{eq:ga_selection_probability}
\end{equation}

在上述表示与评价机制基础上，遗传算法的整体流程可以概括为：首先随机生成初始种群。随后计算种群中各个体的适应度，并依据适应度进行父代选择。在此基础上通过交叉与变异产生新个体，形成下一代种群。最后不断重复上述过程，直至满足最大迭代次数、适应度收敛或达到目标阈值等终止条件。其整体迭代框架如图~\ref{fig:ga_pipeline} 所示，算法从创建初始种群开始，经过适应度评估、选择、交叉和变异不断更新种群，并在满足终止条件后输出当前最优个体。由此可见，遗传算法本质上是一个“编码—评估—选择—交叉—变异”的循环优化过程，其搜索能力来源于种群层面的并行探索以及父代与子代之间的信息传递。

\begin{figure}[!t]
  \centering
  \includegraphics[height=0.5\textheight]{figures/ch2/ga_pipeline}
  \caption{遗传算法的流程图。}
  \label{fig:ga_pipeline}
\end{figure}

在整体流程中，选择（Selection）操作负责依据适应度从当前种群中抽取父代个体。其目的并非简单复制当前最优个体，而是在提高高质量个体生存概率的同时，保留一定程度的群体多样性，从而避免搜索过程过早收敛到局部最优。经典遗传算法中常见的选择方式包括比例选择、排序选择与锦标赛选择等。经过选择后，适应度较高的个体会在统计意义上获得更高的存活与扩增机会，从而为后续的交叉与变异提供更优的遗传材料。

交叉（Crossover）是遗传算法区别于一般随机搜索算法的关键机制之一。其基本思想是：从两个父代染色体中交换部分基因片段，从而生成新的子代个体。对于长度为 $L$ 的二进制串，在最简单的单点交叉中，设随机交叉位置为 $k$，则两个父代可表示为
\begin{equation}
\begin{aligned}
\text{x}_a^{(t)} &= (x_1, \dots, x_k, x_{k+1}, \dots, x_L) \\
\text{x}_b^{(t)} &= (y_1, \dots, y_k, y_{k+1}, \dots, y_L)
\end{aligned}
\label{eq:ga_parents}
\end{equation}
相应地，可产生两个子代
\begin{equation}
\begin{aligned}
\text{z}_1 &= (x_1, \dots, x_k, y_{k+1}, \dots, y_L) \\
\text{z}_2 &= (y_1, \dots, y_k, x_{k+1}, \dots, x_L)
\end{aligned}
\label{eq:ga_offspring}
\end{equation}
交叉的作用并不是简单扰动已有个体，而是通过重组不同父代中的有利局部结构，形成新的候选解组合。换言之，交叉操作承担了组合已有优良结构的作用，是遗传算法实现全局搜索与结构重组的核心机制之一。

变异（Mutation）用于在较低概率下对个体的局部基因进行随机修改，以补充交叉操作难以提供的新信息。对于二进制编码，最常见的变异方式是按位翻转，即对某一位基因执行
\begin{equation}
x_j \leftarrow 1-x_j
\label{eq:ga_mutation}
\end{equation}
若每个位独立以概率 $p_m$ 发生变异，则变异操作能够不断向种群注入新的局部模式，从而维持搜索空间的可达性与群体多样性。经典遗传算法通常将交叉作为主要的结构重组机制，将变异作为防止种群退化与避免早熟收敛的辅助机制。图~\ref{fig:ga_crossover_and_mutate} 给出了遗传算法中交叉与变异操作的示意：其中交叉通过交换两个父代个体的局部基因片段形成新个体，而变异则通过对局部基因位的随机翻转引入新的搜索方向。

\begin{figure}[htbp]
  \centering
  \includegraphics[width=0.6\textwidth,height=0.4\textheight]{figures/ch2/ga_crossover_and_mutate}
  \caption{遗传算法中交叉与变异操作图。}
  \label{fig:ga_crossover_and_mutate}
\end{figure}

在完成选择、交叉与变异之后，遗传算法将生成的新个体组成下一代种群，并重复上述过程，形成如下迭代优化链路：
\begin{equation}
P^{(t)}
\rightarrow \mathrm{Selection}
\rightarrow \mathrm{Crossover}
\rightarrow \mathrm{Mutation}
\rightarrow P^{(t+1)}
\label{eq:ga_evolution_chain}
\end{equation}
随着迭代进行，群体中高适应度个体的比例通常逐步提高，搜索过程由此呈现出从随机探索到定向优化的演化趋势。标准遗传算法的终止条件通常包括达到最大迭代次数、最优个体长期不再改进，或适应度达到预设阈值等。

在二进制编码下，模式可记为
\begin{equation}
H=(h_1, h_2, \dots, h_L), \qquad h_i \in \{0,1,*\}
\label{eq:ga_schema}
\end{equation}
其中符号 $*$ 表示该位可以取任意值，因此一个模式实际上对应解空间中一组具有相似结构的染色体。设 $m(H,t)$ 表示第 $t$ 代种群中属于模式 $H$ 的个体数量，$\bar{f}(H,t)$ 表示该模式实例的平均适应度，$\bar{f}(t)$ 表示整个人群的平均适应度。再记模式的定义长度为 $\delta(H)$，固定位置个数为 $o(H)$，交叉概率为 $p_c$，变异概率为 $p_m$，则经典模式定理给出下一代中该模式期望复制数量的下界：
\begin{equation}
E[m(H,t+1)] \ge
m(H,t)\frac{\bar{f}(H,t)}{\bar{f}(t)}
\left(
1-p_c\frac{\delta(H)}{L-1}-o(H)p_m
\right)
\label{eq:ga_schema_theorem}
\end{equation}
这一结果表明：当某一低阶、短定义长度且平均适应度高于群体平均水平的模式出现时，它在选择作用下更容易得到扩增，而在交叉和变异中的被破坏概率相对较低，因此会在后续代际中以更高概率持续传播。模式定理为遗传算法为何能够通过局部优良结构的保留与重组实现逐步优化，提供了经典的理论解释。

从算法层面看，遗传算法的优势主要体现在三个方面。首先是遗传算法直接作用于编码个体而非问题本身，因此具有较强的表示灵活性。其次是算法通过种群并行搜索减少了对初始点的敏感性，在复杂搜索空间中更容易进行全局探索。最后是选择、交叉与变异的分工使其能够在保留已有优良结构和探索未知候选区域之间取得平衡。与此同时，遗传算法的性能也显著依赖于编码方式、适应度设计、交叉变异概率以及种群更新策略，不合理的参数设置可能导致搜索效率下降或早熟收敛。

综上所述，遗传算法通过“编码—评估—选择—交叉—变异”的统一框架实现对复杂解空间的迭代搜索。这一思想对于候选解具有组合结构、目标函数难以显式求导、评价标准由多项因素共同构成的问题具有较好的适用性。因此，在后续场景生成任务中，若将单个场景方案视为个体，并将场景质量、多样性与任务匹配性等指标共同纳入评价过程，则可自然借鉴遗传算法中的选择、交叉与变异机制，构建面向复杂场景优化的进化式搜索框架。


\section{神经隐式场基础}
\ensection{Neural Implicit Field Foundations}

\subsection{神经辐射场}
\ensubsection{Neural Radiance Fields}
\label{sec:nerf}

神经辐射场（Neural Radiance Fields\cite{mildenhall2020nerf}, NeRF）是一种基于连续隐式函数的三维场景表示算法，其核心思想是利用多视角图像及其相机位姿，通过神经网络拟合三维空间坐标与观测视角到场景体密度及辐射度的映射关系。相比于传统的点云、体素或三角网格等显式表示，NeRF 不直接在离散空间单元上存储场景信息，而是通过神经网络对整个三维空间中的辐射度与密度分布进行连续编码，这种隐式表征特性使其在捕捉复杂几何细节与视点相关的光学效应方面具有显著优势，从而能够实现极高保真度的新视角图像合成。

\begin{figure}[htbp]
    \centering
    \includegraphics[width=0.95\linewidth]{figures/ch2/nerf.png}
    \caption{NeRF 算法框架图\cite{mildenhall2020nerf}。}
    \label{fig:nerf_framework}
\end{figure}

如图~\ref{fig:nerf_framework} 所示，NeRF 将场景表示为一个连续的五维函数。对于空间中任意一点的位置 $\text{x}=(x,y,z)$ 以及对应的观察方向 $\text{d}$，神经网络输出该点的颜色 $\text{c}=(r,g,b)$ 与体积密度 $\sigma$，即
\begin{equation}
F_{\Theta}:(\text{x},\text{d})\rightarrow(\text{c},\sigma),
\label{eq:nerf_function}
\end{equation}
其中，$\Theta$ 表示网络参数。体积密度 $\sigma$ 用于描述空间位置对光线的遮挡或吸收能力，颜色 $\text{c}$ 则依赖于观察方向，用于刻画视角相关外观。

为了增强神经网络对高频几何细节与纹理变化的表达能力，NeRF 在输入端引入位置编码（Positional Encoding）机制，将原始坐标映射到高维频域空间。对于一维标量 $p$，其编码形式可写为
\begin{equation}
\gamma(p)=\left(
\sin(2^{0}\pi p),\cos(2^{0}\pi p),
\sin(2^{1}\pi p),\cos(2^{1}\pi p),
\ldots,
\sin(2^{L-1}\pi p),\cos(2^{L-1}\pi p)
\right),
\label{eq:nerf_pe}
\end{equation}
其中，$L$ 表示频率层数。实际应用时，该映射分别施加在位置向量 $\text{x}$ 和方向向量 $\text{d}$ 的各个分量上。经过位置编码后，网络能够更好地拟合场景中的高频结构，从而提高渲染结果的细节保真度。

在获得辐射场函数后，场景的新视角合成可转化为沿光线的体渲染（Volume Rendering）过程。设一条从相机光心 $\text{o}$ 出发、方向为 $\text{d}$ 的光线为
\begin{equation}
\text{r}(t)=\text{o}+t\text{d},
\label{eq:nerf_ray}
\end{equation}
其中，$t$ 表示沿光线方向的深度参数。则该光线对应的像素颜色可表示为
\begin{equation}
C(\text{r})=\int_{t_n}^{t_f}T(t)\,\sigma(\text{r}(t))\,\text{c}(\text{r}(t),\text{d})\,dt,
\label{eq:nerf_render_continuous}
\end{equation}
其中，$[t_n,t_f]$ 为采样范围，$T(t)$ 表示光线传播到 $t$ 处时未被前方介质遮挡的透射率，其表达式为
\begin{equation}
T(t)=\exp\left(-\int_{t_n}^{t}\sigma(\text{r}(s))\,ds\right).
\label{eq:nerf_transmittance}
\end{equation}
由于连续积分难以直接求解，NeRF 采用离散采样对其进行近似。若在光线上采样得到 $N$ 个点，则像素颜色可写为
\begin{equation}
\hat{C}(\text{r})=\sum_{i=1}^{N}T_i\alpha_i\text{c}_i,
\label{eq:nerf_render_discrete}
\end{equation}
其中
\begin{equation}
\alpha_i = 1-\exp(-\sigma_i\delta_i), \qquad
T_i = \exp\left(-\sum_{j=1}^{i-1}\sigma_j\delta_j\right),
\label{eq:nerf_alpha_t}
\end{equation}
$\sigma_i$ 和 $\text{c}_i$ 分别表示第 $i$ 个采样点的体积密度与颜色，$\delta_i$ 表示相邻采样点之间的距离间隔。由此，NeRF 将图像生成过程转化为沿光线的可微分颜色累积过程，从而能够通过图像监督反向优化场景表示网络。

在训练过程中，NeRF 以真实图像像素颜色作为监督信号，通过最小化预测颜色与图像颜色之间的误差来更新网络参数。对于一批光线集合 $\text{R}$，其损失函数通常写为
\begin{equation}
\text{Loss}=\sum_{\text{r}\in\text{R}}\left\|\hat{C}(\text{r})-C(\text{r})\right\|_2^2,
\label{eq:nerf_loss}
\end{equation}
其中，$C(\text{r})$ 为真实像素颜色，$\hat{C}(\text{r})$ 为渲染得到的预测颜色。为提高采样效率，NeRF 进一步引入分层采样（Hierarchical Sampling）策略，即先利用粗网络进行稀疏采样，再依据粗网络输出的权重分布进行重要性采样，并交由精细网络完成更准确的渲染。该策略能够将更多计算资源分配到对成像更重要的区域，从而在控制计算开销的同时提升渲染质量。


\subsection{神经隐式表面重建}
\ensubsection{NeuS}
\label{sec:neus}

NeuS（Learning Neural Implicit Surfaces by Volume Rendering\cite{wang2021neus}，NeuS）是一种面向多视图重建的神经隐式表面表示算法，其目标是在仅使用多视图图像监督的条件下，学习具有高几何精度的三维隐式表面。针对 NeRF 在几何提取方面的局限性，NeuS 进一步引入有符号距离函数（Signed Distance Function\cite{frisken2000adf}, SDF）作为几何表示，从而将体渲染的优化稳定性与隐式表面的几何精确性结合起来，实现更高精度的几何重建结果。

\begin{figure}[htbp]
    \centering
    \includegraphics[width=0.78\linewidth]{figures/ch2/surface_and_volume_rendering}
    \caption{表面渲染与体渲染的对比图\cite{wang2021neus}。}
    \label{fig:surface_and_volume_rendering}
\end{figure}

NeuS 用一个连续函数 $f_{\theta}(\text{x})$ 表示场景的有符号距离场，其中 $\text{x}\in\mathbb{R}^{3}$ 为空间中的三维点，$\theta$ 为网络参数。对于任意空间点 $\text{x}$，$f_{\theta}(\text{x})$ 表示该点到目标表面的带符号距离：当点位于物体外部时函数值为正，位于内部时函数值为负，而目标表面则由零等值面定义，即
\begin{equation}
S=\{\text{x}\mid f_{\theta}(\text{x})=0\}.
\label{eq:neus_surface}
\end{equation}
在颜色场上，NeuS 继承了 NeRF 的设计，引入颜色场 $c(\text{x},\text{v})$，用于描述空间点 $\text{x}$ 在观察方向 $\text{v}$ 下的颜色。通过这种设计，NeuS 实现了场景几何与颜色的解耦。

在解耦设计下，问题的难点来到如何将 SDF 与图像像素建立可微联系。早期基于表面渲染的算法通常仅考虑每条射线与隐式表面的单次交点，因此梯度主要集中在局部交点附近，当场景中存在严重自遮挡、孔洞或深度突变时，优化过程容易陷入局部最优。相比之下，体渲染算法沿整条光线对多个采样点进行颜色累积，能够为不同深度位置同时提供梯度信号，因此在复杂结构下更具优化稳定性。如图~\ref{fig:surface_and_volume_rendering} 所示，NeuS 的核心思想正是将体渲染的稳定优化优势引入 SDF 表示，从而兼顾隐式表面的几何精确性与训练鲁棒性。

为此，NeuS 首先基于 SDF 构造与表面对应的密度分布。设 $\Phi_s(x)$ 为带参数的 Sigmoid 函数，
\begin{equation}
\Phi_s(x)=\frac{1}{1+e^{-sx}},
\label{eq:neus_sigmoid}
\end{equation}
其中，$s$ 为控制分布尖锐程度的参数。由于 $\Phi_s(x)$ 在零点附近变化最为显著，因此可利用其对 SDF 的响应来刻画表面邻域，并进一步构造体渲染所需的权重分布。

设相机光心为 $\text{o}$，观察方向为 $\text{v}$，则对应光线可写为
\begin{equation}
\text{p}(t)=\text{o}+t\text{v},\quad t\geq 0.
\label{eq:neus_ray}
\end{equation}
NeuS 将该光线对应的像素颜色表示为沿光线的积分累积形式：
\begin{equation}
C(\text{o},\text{v})=\int_{0}^{+\infty} w(t)\, c(\text{p}(t),\text{v})\,dt,
\label{eq:neus_render}
\end{equation}
其中，$w(t)$ 表示深度 $t$ 处采样点对最终像素颜色的贡献权重。与 NeRF 不同，NeuS 的关键在于根据 SDF 重新定义权重函数，使其在表面位置附近取得最大响应，并尽可能减小体渲染带来的几何偏差。

从表面重建角度看，理想的权重函数应满足两个基本要求。其一是无偏性，即当光线穿过真实表面时，权重峰值应与表面交点对齐。其二是遮挡感知性，即当前后存在多个潜在表面时，距离相机更近的表面应具有更大的颜色贡献。若直接将由 SDF 导出的密度代入传统体渲染公式，虽然能够反映遮挡累积，但会导致权重峰值偏离真实表面位置，从而引入几何偏差。如图~\ref{fig:neus_and_naive_solution_bias} 所示，朴素方案的权重峰值出现在真实表面之前，而 NeuS 通过重新构造不透明度密度，使权重函数在一阶近似下能够更好地与表面位置对齐。

\begin{figure}[htbp]
    \centering
    \includegraphics[width=0.7\linewidth]{figures/ch2/neus_and_naive_solution_bias}
    \caption{朴素体渲染方案与 NeuS 无偏体渲染对比图\cite{wang2021neus}。}
    \label{fig:neus_and_naive_solution_bias}
\end{figure}

为实现这一目标，NeuS 保持体渲染中“透射率--密度”乘积的形式，并重新定义不透明度密度 $\rho(t)$：
\begin{equation}
w(t)=T(t)\rho(t),
\label{eq:neus_weight}
\end{equation}
其中透射率 $T(t)$ 定义为
\begin{equation}
T(t)=\exp\left(-\int_{0}^{t}\rho(u)\,du\right).
\label{eq:neus_transmittance}
\end{equation}
在此基础上，NeuS 给出
\begin{equation}
\rho(t)=\max\left(
-\frac{\dfrac{d}{dt}\Phi_s(f_{\theta}(\text{p}(t)))}{\Phi_s(f_{\theta}(\text{p}(t)))},\,0
\right).
\label{eq:neus_opacity_density}
\end{equation}
这种构造方式并非直接使用 SDF 诱导密度，而是通过对 Sigmoid 响应的变化率进行重新定义，使体渲染权重能够在保持遮挡感知性的同时，更准确地聚集在零等值面附近。

在实际计算中，NeuS 采用离散采样对渲染过程进行近似。设沿光线采样得到 $n$ 个点，则像素颜色的离散形式可写为
\begin{equation}
\hat{C}=\sum_{i=1}^{n}T_i\alpha_i\text{c}_i,
\label{eq:neus_discrete_render}
\end{equation}
其中，$\text{c}_i$ 表示第 $i$ 个采样点的颜色，$T_i$ 为累积透射率，$\alpha_i$ 为对应区间的不透明度。NeuS 进一步将 $\alpha_i$ 表达为
\begin{equation}
\alpha_i=
\max\left(
\frac{\Phi_s(f_{\theta}(\text{p}(t_i)))-\Phi_s(f_{\theta}(\text{p}(t_{i+1})))}
{\Phi_s(f_{\theta}(\text{p}(t_i)))},\,0
\right).
\label{eq:neus_alpha}
\end{equation}
由此可见，NeuS 在形式上仍保留了与 NeRF 相似的 $\alpha$ 合成过程，但这些权重由 SDF 经特殊构造后得到，因此能够更好地服务于隐式表面重建。

在训练阶段，NeuS 主要利用渲染颜色与真实图像像素颜色之间的误差来监督网络参数，并引入 Eikonal 正则项约束距离场满足单位梯度性质。其总体损失可写为
\begin{equation}
\text{Loss}
=
\text{Loss}_{\mathrm{color}}
+
\lambda \text{Loss}_{\mathrm{reg}},
\label{eq:neus_total_loss}
\end{equation}
其中，$\text{Loss}_{\mathrm{color}}$ 表示颜色重建误差，$\text{Loss}_{\mathrm{reg}}$ 表示 Eikonal 正则项，$\lambda$ 为权重系数。该正则化能够鼓励网络学习更加平滑且符合 SDF 定义的隐式表面，从而提升几何重建的稳定性。



\section{评价指标}
\ensection{Evaluation Metrics}
\label{sec:evaluation_metrics}


\subsection{场景生成评价指标}
\ensubsection{Scene Generation Metrics}
\noindent\textbf{（1）CLIP分数}

CLIP 分数\cite{hessel2021clipscore}（CLIP Score）是一种用于衡量图像与文本语义一致性的评价指标，常用于文本引导图像生成、跨模态检索以及多模态生成任务中。该指标基于 CLIP（Contrastive Language-Image Pre-training\cite{radford2021learning}，CLIP）模型提取图像与文本的特征表示，并通过计算二者在共享语义空间中的相似度来评价生成结果与文本描述之间的匹配程度。

设输入文本为 $T$，对应图像为 $I$，CLIP 图像编码器与文本编码器分别记为 $f_I$ 和 $f_T$。首先，将图像与文本映射到统一的特征空间中，并对特征向量进行归一化处理，得到
\begin{equation}
\text{v}_I = \frac{f_I(I)}{\|f_I(I)\|_2}, \qquad
\text{v}_T = \frac{f_T(T)}{\|f_T(T)\|_2} .
\label{eq:clip_feature}
\end{equation}

在此基础上，CLIP Score 通常定义为图像特征与文本特征之间的余弦相似度，即
\begin{equation}
\mathrm{CLIPScore}(I,T) = \text{v}_I^\top \text{v}_T .
\label{eq:clipscore}
\end{equation}

由于余弦相似度反映了两个特征向量在语义空间中的方向一致性，因此 CLIP Score 越高，表示图像内容与文本描述在语义上越一致。反之，则表明生成结果与文本条件之间存在较大偏差。在文本驱动的三维场景生成任务中，可以通过渲染生成场景图像，并利用 CLIP 模型计算渲染图像与输入文本之间的相似度，从而评价生成场景在语义对齐方面的性能。相较于仅依赖图像重建误差的评价方式，CLIP Score 更强调高层语义匹配能力，因此适合用于评价文本条件下的场景生成质量。

\noindent\textbf{（2）关键物品覆盖率}

在场景生成任务中，尤其是在文本描述中包含若干明确指定的核心对象时，若生成结果遗漏其中的重要物体，即使整体风格或场景类别与文本条件基本一致，也难以认为其较好地完成了内容生成任务。因此，关键物品覆盖率（Key Object Coverage Ratio）可作为衡量生成结果内容完整性的重要指标，通过统计文本中要求出现的对象在生成结果中是否被正确体现，来衡量模型对输入条件的遵循程度\cite{lin2024instructscene,bai2025freescene}。

设标准关键物品集合为 $\text{O}_{\mathrm{req}}$，其大小为 $N_{\mathrm{required}}$。生成结果中与该集合成功匹配的关键物品数量为 $N_{\mathrm{match}}$，则关键物品覆盖率定义为
\begin{equation}
\mathrm{CoverageRatio}=
\frac{N_{\mathrm{match}}}{N_{\mathrm{required}}}.
\label{eq:key_object_coverage}
\end{equation}

其中，$N_{\mathrm{required}}$ 表示参考物品集合中要求包含的关键物品总数，$N_{\mathrm{match}}$ 表示生成结果中成功出现并与参考集合语义对应的关键物品数量。该指标的取值范围通常为 $[0,1]$。其值越高，表明生成结果对目标关键物品的覆盖越完整。反之，则表明生成过程中存在物体遗漏或内容缺失现象。

关键物品覆盖率能够较为直接地反映模型在对象层面对输入条件的满足能力，尤其适合用于评价生成结果中“应出现的物体是否被生成”这一问题。在文本驱动的三维场景生成任务中，该指标可用于分析模型对核心对象、功能性物体或约束指定物品的保留能力，从而补充整体语义相似度指标对局部对象完整性刻画不足的问题。

\noindent\textbf{（3）人工评分}

人工评分是一种由人工评审者根据预设评价维度对生成场景进行主观评分的评价方式。作为一种主观评价指标，人工评分能够从用户感知角度对生成结果进行综合评价，尤其适合补充客观指标在布局合理性、视觉真实性与整体可用性刻画上的不足。对于文本驱动的三维场景生成任务，特别是面向化学实验的专业场景，人工评价还可进一步辅助判断安全设施是否完整、实验布置是否符合常识，以及场景是否满足实际使用预期，具有重要作用。



\subsection{三维重建评价指标}
\ensubsection{3D Reconstruction Metrics}
\noindent\textbf{（1）峰值信噪比}

峰值信噪比（Peak Signal-to-Noise Ratio\cite{huynhthu2008psnr}，PSNR）是图像重建、图像压缩与图像生成任务中常用的客观评价指标，用于衡量重建结果相对于参考图像的失真程度。PSNR 本质上是基于原始图像与重建图像之间的均方误差（Mean Squared Error，MSE）定义的。一般而言，PSNR 值越高，表明重建图像与参考图像越接近，图像失真越小，重建质量越好。

设原始图像为 $I$，重建图像为 $\hat{I}$，图像分辨率为 $m \times n$，则二者之间的均方误差可表示为
\begin{equation}
\mathrm{MSE} = \frac{1}{mn} \sum_{i=0}^{m-1} \sum_{j=0}^{n-1} \left( I(i,j) - \hat{I}(i,j) \right)^2 ,
\label{eq:psnr_mse}
\end{equation}
其中，$I(i,j)$ 和 $\hat{I}(i,j)$ 分别表示原始图像与重建图像在像素坐标 $(i,j)$ 处的像素值。

在得到均方误差后，PSNR 可进一步定义为
\begin{equation}
\mathrm{PSNR} = 20 \log_{10}\left( \frac{\mathrm{MAX}_I}{\sqrt{\mathrm{MSE}}} \right) ,
\label{eq:psnr}
\end{equation}
其中，$\mathrm{MAX}_I$ 表示图像像素的最大可能取值。对于 8 bit 图像，通常有 $\mathrm{MAX}_I = 255$。

从公式可以看出，PSNR 与 MSE 呈单调负相关关系：当重建图像与原始图像越接近时，$\mathrm{MSE}$ 越小，对应的 $\mathrm{PSNR}$ 越大。因此，PSNR 能够直观反映重建结果在像素层面的保真程度。在三维重建任务中，通常将模型渲染图像与真实采集图像进行逐像素比较，并通过计算 PSNR 来评价模型在外观重建方面的准确性。

\noindent\textbf{（2）结构相似性}

结构相似性（Structural Similarity Index Measure\cite{wang2004ssim}，SSIM）是一种用于衡量两幅图像在结构信息层面相似程度的图像质量评价指标，常用于图像重建、图像生成与三维重建等任务中。与 PSNR 基于逐像素误差不同，SSIM 更关注图像在亮度、对比度和结构三方面的整体一致性，因此能够在一定程度上更符合人类视觉系统对图像质量的主观感知。

设原始图像为 $I$，重建图像为 $\hat{I}$，分别记其局部区域的均值为 $\mu_I$ 和 $\mu_{\hat{I}}$，标准差为 $\sigma_I$ 和 $\sigma_{\hat{I}}$，协方差为 $\sigma_{I\hat{I}}$。则 SSIM 可定义为
\begin{equation}
\mathrm{SSIM}(I,\hat{I}) =
\frac{(2\mu_I\mu_{\hat{I}} + C_1)(2\sigma_{I\hat{I}} + C_2)}
{(\mu_I^2 + \mu_{\hat{I}}^2 + C_1)(\sigma_I^2 + \sigma_{\hat{I}}^2 + C_2)} ,
\label{eq:ssim}
\end{equation}
其中，$C_1$ 和 $C_2$ 为两个用于避免分母过小而导致数值不稳定的常数，通常取为
\begin{equation}
C_1 = (K_1 L)^2,\qquad C_2 = (K_2 L)^2,
\label{eq:ssim_constants}
\end{equation}
其中，$L$ 表示图像像素值的动态范围，$K_1$ 和 $K_2$ 为经验常数，常取 $K_1=0.01$、$K_2=0.03$。

从公式可以看出，SSIM 并非直接计算像素差值，而是分别从亮度一致性、对比度一致性和结构一致性三个角度对图像进行综合评价。其取值范围通常在 $[0,1]$ 之间，数值越大表示两幅图像在结构信息上越接近，重建结果越好。当 $\mathrm{SSIM}=1$ 时，表示两幅图像完全一致。

在三维重建任务中，SSIM 通常用于比较模型渲染图像与真实采集图像之间的相似程度，以评价重建结果在局部结构保持、边缘质量以及整体视觉一致性方面的表现。相较于 PSNR，SSIM 对亮度偏移和局部结构变化更敏感，因此能够从结构保真的角度补充反映模型的重建质量。由于其兼顾数学定义与感知一致性，SSIM 已成为图像重建和新视角合成任务中的常用评价指标之一。

\noindent\textbf{（3）学习感知图像块相似度}

学习感知图像块相似度（Learned Perceptual Image Patch Similarity\cite{zhang2018lpips}，LPIPS）是一种基于深度特征的感知质量评价指标，用于衡量两幅图像在人类视觉感知层面的差异。与 PSNR 和 SSIM 主要依赖像素级误差或局部统计量不同，LPIPS 借助预训练卷积神经网络提取图像的多层特征表示，从而能够更好地反映图像在纹理、结构和语义层面的感知差异。

设两幅输入图像分别为 $x$ 和 $y$，通过预训练网络第 $l$ 层提取得到的特征图分别为 $F_l(x)$ 和 $F_l(y)$。对特征进行通道归一化后，LPIPS 可表示为
\begin{equation}
\mathrm{LPIPS}(x,y) =
\sum_l \frac{1}{H_l W_l}
\sum_{h,w}
\left\|
w_l \odot
\left( \hat{F}_l(x)_{hw} - \hat{F}_l(y)_{hw} \right)
\right\|_2^2 ,
\label{eq:lpips}
\end{equation}
其中，$\hat{F}_l$ 表示归一化后的特征，$H_l$ 和 $W_l$ 分别表示第 $l$ 层特征图的高和宽，$w_l$ 表示该层特征的线性加权系数，$\odot$ 表示逐通道乘法。

LPIPS 的数值越小，表示两幅图像在感知层面越相似，重建结果越接近真实图像。在图像生成与三维重建任务中，LPIPS 常用于衡量模型输出与参考图像之间的感知差异，尤其适合评估纹理细节和视觉真实感。

相较于传统的像素级评价指标，LPIPS 更符合人类主观感知结果，因此在重建质量评价中具有重要作用。然而，由于其依赖预训练特征网络，不同网络结构和训练设置可能会对评价结果产生一定影响。因此，LPIPS 一般作为感知质量的补充指标，与 PSNR 和 SSIM 共同用于综合评估模型性能。

\noindent\textbf{（4）Chamfer距离}

Chamfer 距离（Chamfer Distance\cite{fan2017chamfer}，CD）是三维点集与几何重建任务中常用的形状相似性评价指标，用于衡量预测点集与真实点集之间的双向距离差异。在三维重建中，通常将预测网格表面与真实网格表面均匀采样为点集，再通过 Chamfer 距离评价两者在几何形状上的接近程度。

设真实点集为 $P=\{p_i\}_{i=1}^{N}$，预测点集为 $Q=\{q_j\}_{j=1}^{M}$，则 Chamfer 距离可定义为
\begin{equation}
\mathrm{CD}(P,Q) =
\frac{1}{N}\sum_{p \in P}\min_{q \in Q}\|p-q\|_2^2
+
\frac{1}{M}\sum_{q \in Q}\min_{p \in P}\|q-p\|_2^2 ,
\label{eq:cd}
\end{equation}
其中，$\|\cdot\|_2$ 表示欧氏距离。

上述公式从两个方向度量点集之间的接近程度：第一项表示真实点集中的每个点到预测点集最近点的平均距离，用于衡量预测结果对真实形状的覆盖情况。第二项表示预测点集中的每个点到真实点集最近点的平均距离，用于衡量预测结果中是否存在偏离真实表面的冗余部分。通过双向约束，Chamfer 距离能够综合刻画预测几何与真实几何之间的整体差异。

Chamfer 距离的数值越小，表示预测几何与真实几何越接近，重建质量越高。相较于基于图像的评价指标，Chamfer 距离直接作用于三维几何本身，因此更适合评价模型的表面恢复能力、细节保留能力以及整体形状精度。在三维重建实验中，Chamfer 距离常作为定量分析几何重建质量的重要指标之一。

\subsection{大语言模型生成质量评价指标}
\ensubsection{LLM Generation Metrics}

在大语言模型从自然语言指令到结构化输出的生成任务中，评价指标不仅需要衡量输出结果在语义层面是否正确，还需要关注其是否满足下游执行系统的接口约束。因此，相关工作\cite{wei2025plangenllms,yu2018spider,zhong2020semantic}通常从整体匹配、参数匹配以及可执行性三个角度对生成结果进行评估。对于具有标准参考答案的结构化规划任务，整体精确匹配用于衡量预测结果与标准结果是否完全一致。对于包含多个字段或参数的输出，还可采用参数级匹配指标评价局部正确性。当预测结果需要被执行器或仿真系统直接调用时，还需要进一步考察其可执行性。

设测试集包含 $N$ 个样本，第 $i$ 个样本的标准动作序列记为 $Y_i$，预测动作序列记为 $\hat{Y}_i$。

\textbf{（1）动作正确率（Action Accuracy, AA）。}
若第 $i$ 个样本的预测动作序列与标准答案完全一致，则记
\[
e_i =
\begin{cases}
1, & \hat{Y}_i = Y_i,\\
0, & \text{otherwise}.
\end{cases}
\]
则动作精确匹配率定义为
\begin{equation}
\mathrm{AA}=\frac{1}{N}\sum_{i=1}^{N} e_i .
\label{eq:action_accuracy}
\end{equation}
该指标反映模型能否完整生成与参考答案一致的结构化动作序列。

\textbf{（2）参数准确率。（Parameter Accuracy, PA）}
设第 $i$ 个样本的标准动作序列中共有 $m_i$ 个待评估参数，其中预测准确的参数个数为 $c_i$，则该样本的参数匹配率为
\[
p_i=\frac{c_i}{m_i}.
\]
进一步地，全测试集上的参数级匹配率定义为
\begin{equation}
\mathrm{PA}=\frac{1}{N}\sum_{i=1}^{N} p_i .
\label{eq:param_match_rate}
\end{equation}
该指标用于衡量模型对操作对象、目标容器、体积、质量和持续时间等关键信息的预测准确程度。

\textbf{（3）成功率（Success Rate, SR）。}
若第 $i$ 个样本的预测结果能够通过静态接口检查，则记
\[
r_i=
\begin{cases}
1, & \hat{Y}_i \text{ is executable},\\
0, & \text{otherwise}.
\end{cases}
\]
则可执行率定义为
\begin{equation}
\mathrm{SR}=\frac{1}{N}\sum_{i=1}^{N} r_i .
\label{eq:executability_rate}
\end{equation}
其中，可执行是指该动作序列满足下游执行器要求的动作类型、必需参数和格式规范。该指标反映生成结果能否作为控制模块的有效输入。

综上所述，动作准确率、参数准确率与成功率分别从整体结构一致性、局部参数正确性以及下游执行可用性三个方面评价模型输出质量。

\section{本章小结}
\ensection{Summary}

本章围绕本文后续研究所涉及的关键技术，系统介绍了相关理论基础，为第三章、第四章和第五章的算法设计与实验分析提供了支撑。首先，阐述了 Transformer 的基本原理，以及上下文学习、思维链和 Zero-shot-CoT 等提示工程算法，表明了大语言模型完成结构化理解与推理生成的内在机制。其次，介绍了遗传算法的核心思想与基本流程，包括选择、交叉和变异等关键操作，为后续场景进化优化算法奠定了理论基础。随后，系统分析了 NeRF 与 NeuS 的基本原理，表明了神经隐式表示、体渲染以及基于有符号距离函数的表面重建机制，为后续高保真三维重建算法提供了支撑。最后，总结了本文实验中使用的主要评价指标，为后续实验结果分析提供了统一的评价依据。